\usepackage{natbib}         % Pour la bibliographie
\usepackage{url}            % Pour citer les adresses web
%\usepackage[T1]{fontenc}    % Encodage des accents
%\usepackage{tgbonum}
\usepackage{cmbright}
\usepackage[T1]{fontenc}
\usepackage[utf8]{inputenc} % Lui aussi
\usepackage[french]{babel} % Pour la traduction française
\usepackage{numprint}       % Histoire que les chiffres soient bien

\usepackage{amsmath}        % La base pour les maths
\usepackage{mathrsfs}       % Quelques symboles supplémentaires
\usepackage{amssymb}        % encore des symboles.
\usepackage{amsfonts}       % Des fontes, eg pour \mathbb.
\usepackage{amsthm,bbm}
\usepackage{mathtools}
\usepackage{xcolor}

\usepackage{cancel}
\usepackage{pdfpages}

\usepackage[table]{xcolor} % for coloring cells
\usepackage{array} % for column formatting

%\usepackage[svgnames]{xcolor} % De la couleur

%%% Si jamais vous voulez changer de police: décommentez les trois 
%\usepackage{tgpagella}
%\usepackage{tgadventor}
%\usepackage{inconsolata}

%%% Code highlighting Python
%\usepackage{minted}
%\usemintedstyle{friendly}

\usepackage{graphicx} % inclusion des graphiques
\usepackage{wrapfig}  % Dessins dans le texte.

\usepackage{standalone} % pour compiler tikz indep

\usepackage{tikz}     % Un package pour les dessins (utilisé pour l'environnement {code})
%\usetikzlibrary{external}
%\tikzexternalize[prefix=tikz/]

\usepackage{pgfplots}
\usetikzlibrary{shapes.misc}
\pgfplotsset{compat=1.18}
\usepackage[framemethod=TikZ]{mdframed}
\usepackage{multimedia}
\usepackage{media9}
\usepackage{subcaption}
\usepackage{float}
\usepackage[most]{tcolorbox}
\usepackage{enumerate}
\usetikzlibrary{calc}
%mise en page

\usepackage{geometry}
 \geometry{
 a4paper,
 total={170mm,257mm},
 left=20mm,
 top=20mm,
 right=20mm,
 bottom=20mm
 }
\usepackage{setspace}

\onehalfspacing
%\fontfamily{cmss}\selectfont
\makeatletter
\newcommand\subsubsubsection{\@startsection{paragraph}{4}{\z@}%
            {-2.5ex\@plus -1ex \@minus -.25ex}%
            {1.25ex \@plus .25ex}%
            {\normalfont\normalsize\bfseries}}
\setcounter{secnumdepth}{4} % how many sectioning levels to assign numbers to
\setcounter{tocdepth}{2}    % how many sectioning levels to show in ToC

\newcommand{\comments}[1]{}